% Options for packages loaded elsewhere
% Options for packages loaded elsewhere
\PassOptionsToPackage{unicode}{hyperref}
\PassOptionsToPackage{hyphens}{url}
\PassOptionsToPackage{dvipsnames,svgnames,x11names}{xcolor}
%
\documentclass[
  letterpaper,
  DIV=11,
  numbers=noendperiod]{scrreport}
\usepackage{xcolor}
\usepackage[top=30mm,left=20mm,heightrounded]{geometry}
\usepackage{amsmath,amssymb}
\setcounter{secnumdepth}{5}
\usepackage{iftex}
\ifPDFTeX
  \usepackage[T1]{fontenc}
  \usepackage[utf8]{inputenc}
  \usepackage{textcomp} % provide euro and other symbols
\else % if luatex or xetex
  \usepackage{unicode-math} % this also loads fontspec
  \defaultfontfeatures{Scale=MatchLowercase}
  \defaultfontfeatures[\rmfamily]{Ligatures=TeX,Scale=1}
\fi
\usepackage[]{libertinus}
\ifPDFTeX\else
  % xetex/luatex font selection
\fi
% Use upquote if available, for straight quotes in verbatim environments
\IfFileExists{upquote.sty}{\usepackage{upquote}}{}
\IfFileExists{microtype.sty}{% use microtype if available
  \usepackage[]{microtype}
  \UseMicrotypeSet[protrusion]{basicmath} % disable protrusion for tt fonts
}{}
\makeatletter
\@ifundefined{KOMAClassName}{% if non-KOMA class
  \IfFileExists{parskip.sty}{%
    \usepackage{parskip}
  }{% else
    \setlength{\parindent}{0pt}
    \setlength{\parskip}{6pt plus 2pt minus 1pt}}
}{% if KOMA class
  \KOMAoptions{parskip=half}}
\makeatother
% Make \paragraph and \subparagraph free-standing
\makeatletter
\ifx\paragraph\undefined\else
  \let\oldparagraph\paragraph
  \renewcommand{\paragraph}{
    \@ifstar
      \xxxParagraphStar
      \xxxParagraphNoStar
  }
  \newcommand{\xxxParagraphStar}[1]{\oldparagraph*{#1}\mbox{}}
  \newcommand{\xxxParagraphNoStar}[1]{\oldparagraph{#1}\mbox{}}
\fi
\ifx\subparagraph\undefined\else
  \let\oldsubparagraph\subparagraph
  \renewcommand{\subparagraph}{
    \@ifstar
      \xxxSubParagraphStar
      \xxxSubParagraphNoStar
  }
  \newcommand{\xxxSubParagraphStar}[1]{\oldsubparagraph*{#1}\mbox{}}
  \newcommand{\xxxSubParagraphNoStar}[1]{\oldsubparagraph{#1}\mbox{}}
\fi
\makeatother


\usepackage{longtable,booktabs,array}
\usepackage{calc} % for calculating minipage widths
% Correct order of tables after \paragraph or \subparagraph
\usepackage{etoolbox}
\makeatletter
\patchcmd\longtable{\par}{\if@noskipsec\mbox{}\fi\par}{}{}
\makeatother
% Allow footnotes in longtable head/foot
\IfFileExists{footnotehyper.sty}{\usepackage{footnotehyper}}{\usepackage{footnote}}
\makesavenoteenv{longtable}
\usepackage{graphicx}
\makeatletter
\newsavebox\pandoc@box
\newcommand*\pandocbounded[1]{% scales image to fit in text height/width
  \sbox\pandoc@box{#1}%
  \Gscale@div\@tempa{\textheight}{\dimexpr\ht\pandoc@box+\dp\pandoc@box\relax}%
  \Gscale@div\@tempb{\linewidth}{\wd\pandoc@box}%
  \ifdim\@tempb\p@<\@tempa\p@\let\@tempa\@tempb\fi% select the smaller of both
  \ifdim\@tempa\p@<\p@\scalebox{\@tempa}{\usebox\pandoc@box}%
  \else\usebox{\pandoc@box}%
  \fi%
}
% Set default figure placement to htbp
\def\fps@figure{htbp}
\makeatother


% definitions for citeproc citations
\NewDocumentCommand\citeproctext{}{}
\NewDocumentCommand\citeproc{mm}{%
  \begingroup\def\citeproctext{#2}\cite{#1}\endgroup}
\makeatletter
 % allow citations to break across lines
 \let\@cite@ofmt\@firstofone
 % avoid brackets around text for \cite:
 \def\@biblabel#1{}
 \def\@cite#1#2{{#1\if@tempswa , #2\fi}}
\makeatother
\newlength{\cslhangindent}
\setlength{\cslhangindent}{1.5em}
\newlength{\csllabelwidth}
\setlength{\csllabelwidth}{3em}
\newenvironment{CSLReferences}[2] % #1 hanging-indent, #2 entry-spacing
 {\begin{list}{}{%
  \setlength{\itemindent}{0pt}
  \setlength{\leftmargin}{0pt}
  \setlength{\parsep}{0pt}
  % turn on hanging indent if param 1 is 1
  \ifodd #1
   \setlength{\leftmargin}{\cslhangindent}
   \setlength{\itemindent}{-1\cslhangindent}
  \fi
  % set entry spacing
  \setlength{\itemsep}{#2\baselineskip}}}
 {\end{list}}
\usepackage{calc}
\newcommand{\CSLBlock}[1]{\hfill\break\parbox[t]{\linewidth}{\strut\ignorespaces#1\strut}}
\newcommand{\CSLLeftMargin}[1]{\parbox[t]{\csllabelwidth}{\strut#1\strut}}
\newcommand{\CSLRightInline}[1]{\parbox[t]{\linewidth - \csllabelwidth}{\strut#1\strut}}
\newcommand{\CSLIndent}[1]{\hspace{\cslhangindent}#1}



\setlength{\emergencystretch}{3em} % prevent overfull lines

\providecommand{\tightlist}{%
  \setlength{\itemsep}{0pt}\setlength{\parskip}{0pt}}



 


\KOMAoption{captions}{tableheading}
\makeatletter
\@ifpackageloaded{bookmark}{}{\usepackage{bookmark}}
\makeatother
\makeatletter
\@ifpackageloaded{caption}{}{\usepackage{caption}}
\AtBeginDocument{%
\ifdefined\contentsname
  \renewcommand*\contentsname{Table of contents}
\else
  \newcommand\contentsname{Table of contents}
\fi
\ifdefined\listfigurename
  \renewcommand*\listfigurename{List of Figures}
\else
  \newcommand\listfigurename{List of Figures}
\fi
\ifdefined\listtablename
  \renewcommand*\listtablename{List of Tables}
\else
  \newcommand\listtablename{List of Tables}
\fi
\ifdefined\figurename
  \renewcommand*\figurename{Figure}
\else
  \newcommand\figurename{Figure}
\fi
\ifdefined\tablename
  \renewcommand*\tablename{Table}
\else
  \newcommand\tablename{Table}
\fi
}
\@ifpackageloaded{float}{}{\usepackage{float}}
\floatstyle{ruled}
\@ifundefined{c@chapter}{\newfloat{codelisting}{h}{lop}}{\newfloat{codelisting}{h}{lop}[chapter]}
\floatname{codelisting}{Listing}
\newcommand*\listoflistings{\listof{codelisting}{List of Listings}}
\makeatother
\makeatletter
\makeatother
\makeatletter
\@ifpackageloaded{caption}{}{\usepackage{caption}}
\@ifpackageloaded{subcaption}{}{\usepackage{subcaption}}
\makeatother
\usepackage{bookmark}
\IfFileExists{xurl.sty}{\usepackage{xurl}}{} % add URL line breaks if available
\urlstyle{same}
\hypersetup{
  pdftitle={Final Project Report},
  pdfauthor={Kanika Dhawan , Mansi Yadav ,Prince , Vishal , Rahul Nayak},
  colorlinks=true,
  linkcolor={blue},
  filecolor={Maroon},
  citecolor={Blue},
  urlcolor={Blue},
  pdfcreator={LaTeX via pandoc}}


\title{Final Project Report}
\author{Kanika Dhawan , Mansi Yadav ,Prince , Vishal , Rahul Nayak}
\date{2025-06-16}
\begin{document}
\maketitle

\renewcommand*\contentsname{Table of contents}
{
\hypersetup{linkcolor=}
\setcounter{tocdepth}{2}
\tableofcontents
}
\listoffigures
\listoftables

\bookmarksetup{startatroot}

\chapter*{Preface}\label{preface}
\addcontentsline{toc}{chapter}{Preface}

\markboth{Preface}{Preface}









\bookmarksetup{startatroot}

\chapter*{Abstract}\label{abstract}
\addcontentsline{toc}{chapter}{Abstract}

\markboth{Abstract}{Abstract}

This project focuses on developing an AI-based system to monitor and
predict driver behavior with the goal of improving road safety. By using
video analysis, GPS data, and vehicle parameters such as speed,
acceleration, and braking, the system detects unsafe driving patterns in
real-time. It provides timely feedback to drivers, helping them correct
risky behavior and adopt safer driving habits. The main objective is to
reduce accidents and promote responsible driving through intelligent,
data-driven insights

\section*{Acknowledgment}\label{acknowledgment}
\addcontentsline{toc}{section}{Acknowledgment}

\markright{Acknowledgment}

This project has been a significant learning experience, combining
multiple disciplines such as computer vision, deep learning, and
real-time system integration. I would like to acknowledge the valuable
assistance of AI tools, particularly ChatGPT, which supported the
writing process of this report. The AI was utilized to improve the
clarity, coherence, and technical articulation of complex concepts, as
well as to ensure proper formatting and structure throughout the
documentation.

It is important to emphasize that while generative AI tools contributed
to enhancing the written presentation, all core development work ---
including dataset curation, preprocessing, model training, evaluation,
optimization, and real-time inference --- was independently
conceptualized and implemented by me. The ideas, coding logic,
architecture design, experimentation, and debugging reflect original
effort, critical thinking, and practical application of machine learning
and computer vision knowledge.

Furthermore, all external resources, including datasets, pre-trained
models, or open-source libraries, have been properly cited in the
relevant sections of this report. The final system represents a genuine
and hands-on implementation of a real-time driver monitoring solution,
grounded in both theoretical understanding and practical
experimentation.

Special thanks to the mentors, teammates, and the institution (Sabudh
Foundation) for their guidance and support throughout the course of this
project.

\bookmarksetup{startatroot}

\chapter{Introduction to Project}\label{introduction-to-project}

human mistake is one of the main causes of traffic accidents, road
safety is a big concern worldwide. In order to reduce accidents and
raise driving standards generally, it can be quite important to keep an
eye on driver conduct and give prompt criticism. Artificial Intelligence
(AI) has become a potent instrument in recent years for analyzing
intricate human behaviors and making data-driven judgments in real time.

The goal of this project is to create an AI-powered system that uses GPS
data, video analysis and vehicle characteristics like speed,
acceleration, and braking to forecast and assess driver behavior. The
ultimate goal is to identify unsafe driving habits in real-time and
provide proactive feedback to drivers. This not only helps reduce the
risk of accidents but also encourages the development of safer and more
responsible driving practices

\section{Overview}\label{overview}

This project is focused on using AI to make driving safer and smarter.
Instead of just relying on basic driving data, it brings together a
bunch of information to get a full picture of how the driver is behaving
behind the wheel. Here's how it works:

\begin{enumerate}
\def\labelenumi{\arabic{enumi}.}
\item
  It checks how alert the driver is --- whether they're focused or
  getting distracted.
\item
  It looks for signs of drowsiness, like drooping eyelids or yawning.
\item
  It even watches for tiny facial expressions that can show stress or
  discomfort.
\end{enumerate}

To do this, the system uses:

\begin{enumerate}
\def\labelenumi{\arabic{enumi}.}
\item
  A live camera feed to analyze the driver's face and head movements.
\item
  GPS data to understand driving routes and patterns.
\item
  Vehicle data like speed, acceleration, and braking behavior.
\end{enumerate}

Once all this data is collected, the AI processes it and calculates a
unified safety score that sums up how safely the driver is performing at
any moment. If something seems off, like if the driver looks sleepy or
stressed, the system can give real-time feedback or alerts to help them
correct their behavior.

In short, the goal is to spot risky driving habits early and guide
drivers toward safer, more responsible driving --- using the power of AI
to make roads safer for everyone.

\section{ExistingSystem}\label{existingsystem}

Currently, several systems and technologies aim to improve driver
safety, but most of them have limitations when it comes to analyzing
driver behavior in a detailed and real-time manner. Some of the existing
solutions include:

\begin{enumerate}
\def\labelenumi{\arabic{enumi}.}
\item
  \textbf{Basic Driver Monitoring Systems}: Many modern vehicles come
  with built-in alert systems that beep if the driver doesn't wear a
  seatbelt, drifts from the lane, or brakes suddenly. However, these
  systems usually rely only on sensor data and don't analyze the
  driver's actual behavior or condition.
\item
  \textbf{Drowsiness Detection Systems:} Some cars have simple
  drowsiness detection features based on steering patterns or eyelid
  movement. These systems may not be accurate in all situations and
  often trigger alerts only after the driver is already at risk.
\item
  \textbf{GPS-based Tracking Systems:} Fleet vehicles and taxis often
  use GPS tracking to monitor speed, location, and route history. These
  systems are good for logging data but do not offer real-time insights
  into the driver's physical or emotional state.
\item
  \textbf{Third-party mobile apps:} There are apps that use phone
  sensors to monitor driving speed or phone usage while driving.
  However, they depend heavily on the phone's position and permissions,
  which can affect accuracy.
\end{enumerate}

\textbf{Limitations of the Existing Systems:} 1. Lack of real-time
behavior prediction

\begin{enumerate}
\def\labelenumi{\arabic{enumi}.}
\setcounter{enumi}{1}
\item
  Limited focus on driver's physical and emotional state
\item
  No integration of facial expression or alertness analysis
\end{enumerate}

These limitations show that there's a need for a more intelligent,
multi-input system that can understand a driver's condition in real time
and take action before an accident happens.

\section{Objectives of Project (Must be clearly, precisely defined and
must be covered in the
work)}\label{objectives-of-project-must-be-clearly-precisely-defined-and-must-be-covered-in-the-work}

The primary objective of this project is to build an AI-based system
that can monitor and assess driver behavior in real time to improve road
safety. The system aims to generate a unified driver safety score by
analyzing various inputs such as facial expressions, alertness, and
drowsiness. It uses live video to detect signs of distraction or
fatigue, tracks subtle facial micro-expressions to gauge emotional
states like stress or anger, and monitors eye and head movements to
identify drowsiness. Additionally, the system collects vehicle data such
as speed, acceleration, and braking patterns, along with GPS data to
understand driving routes and behavior in different environments. Based
on this combined analysis, the system provides real-time feedback to the
driver, helping them correct unsafe behavior before it leads to an
accident. The overall goal is to create a proactive, intelligent
solution that encourages safer and more responsible driving habits.

\bookmarksetup{startatroot}

\chapter{Pre-Processing and Exploratory Data
Analysis}\label{pre-processing-and-exploratory-data-analysis}

\subsection{Pre-Processing}\label{pre-processing}

Pre-processing refers to the set of techniques used to clean and prepare
raw data for input into a machine learning model. This includes
operations such as resizing images, normalizing data, handling missing
values, encoding categorical variables, and generating sequences or
features required by the model. Effective pre-processing ensures data
consistency, reduces noise, and improves model performance.

\subsection{Exploratory Data Analysis
(EDA)}\label{exploratory-data-analysis-eda}

Exploratory Data Analysis (EDA) is the process of analyzing and
visualizing datasets to uncover patterns, spot anomalies, test
hypotheses, and check assumptions. It involves techniques such as
plotting distributions, inspecting class balance, and summarizing
statistical properties of the features. EDA provides critical insights
that guide model selection, feature engineering, and validation
strategies. \#\# Dataset Collection

The dataset used in this project comes from multiple sources, including:

\begin{itemize}
\item
  \textbf{State Farm Distracted Driver Detection Dataset}\\
  A publicly available dataset containing real-world images of drivers
  performing various activities (e.g., texting, talking on the phone,
  eating), categorized into 10 labeled classes.
\item
  \textbf{Custom DAD-style Video Dataset}\\
  Manually recorded driver videos with annotations for focus level,
  drowsiness, and emotional microexpressions. These videos were
  segmented and labeled to match the frame rate of the real-time system
  (30 FPS).
\end{itemize}

These datasets provide diverse conditions, angles, and driver
demographics, which help the model generalize better during training and
inference.

\begin{center}\rule{0.5\linewidth}{0.5pt}\end{center}

\section{Data Pre-processing}\label{data-pre-processing}

Effective data pre-processing is crucial to ensure the quality and
consistency of the input data before feeding it to the machine learning
models.

\subsection{Handling Missing Values}\label{handling-missing-values}

Missing values can arise due to face detection failures or inconsistent
landmark extraction in some video frames. To handle this: - Frames with
incomplete landmark data are discarded. - For sequence-based models like
LSTM, only complete 30-frame sequences are retained. - In some cases,
interpolation is used to fill short gaps.

\subsection{Data Size and
Dimensionality}\label{data-size-and-dimensionality}

High frame rates and video-based data can lead to large volumes of data:
- Each video is broken down into 30-frame sequences. - Feature vectors
such as facial keypoints are extracted to reduce dimensionality. -
Efficient storage formats (e.g., \texttt{.npy}, \texttt{.pt}) are used
to manage memory overhead.

Some deep learning algorithms (e.g., LSTM) are more sensitive to long
sequences, so chunking is used to make training more efficient.

\subsection{Scaling of Variables}\label{scaling-of-variables}

Scaling is essential when working with models like LSTM or ResNet: -
Facial keypoints are normalized relative to face size. - Pixel values in
image-based models are scaled between 0 and 1. - This ensures that no
feature disproportionately influences the model.

\subsection{Outliers}\label{outliers}

Outliers, such as blurred frames or extreme head turns, can mislead
training: - Detected using visualization and statistical thresholds. -
Either removed or smoothed using temporal filters (moving average,
median filters).

\subsection{Clustering (if applicable)}\label{clustering-if-applicable}

In the microexpression model, facial embeddings can form clusters
corresponding to different emotional states: - Dimensionality reduction
techniques like PCA or t-SNE were used to visualize clusters. - K-means
was tested to confirm separability, and the majority-voting approach
helped in aggregating weak predictions.

\subsection{Data Transformation}\label{data-transformation}

Several transformations improve learning: - For the attentiveness model,
directional vectors (nose to eyes, etc.) are computed from landmarks. -
Grayscale conversion simplifies emotion detection without losing crucial
features. - Temporal smoothing over sequence windows helps reduce noise.

\begin{center}\rule{0.5\linewidth}{0.5pt}\end{center}

Overall, these pre-processing strategies ensure that the models receive
clean, normalized, and structured input data that enhances their
predictive power and generalization.

\section{Exploratory Data Analysis and
Visualisations}\label{exploratory-data-analysis-and-visualisations}

Exploratory Data Analysis (EDA) plays a crucial role in understanding
the structure, quality, and hidden patterns of the dataset before
applying machine learning models. It helps in making informed decisions
during model design and improves overall system performance.

\subsection{1. Feature Selection}\label{feature-selection}

Feature selection identifies the most informative and relevant variables
for modeling: - For attentiveness detection, directional vectors (from
nose to eyes, etc.) are prioritized. - For drowsiness detection,
features like EAR and blink rate are retained. - For emotion detection,
spatial patterns in grayscale pixel intensity are emphasized.

Redundant or non-contributing features are removed to reduce overfitting
and improve model interpretability.

\subsection{2. Dimensionality Reduction}\label{dimensionality-reduction}

When dealing with high-dimensional data such as facial landmarks or
CNN-generated embeddings: - \textbf{Principal Component Analysis (PCA)}
and \textbf{t-SNE} are employed to reduce feature dimensions while
preserving variance. - Dimensionality reduction improves model training
time and helps in visualizing data clusters (e.g., emotion categories).

These techniques also help in validating whether the features separate
well across different classes, which is critical for classifier
effectiveness.

\begin{center}\rule{0.5\linewidth}{0.5pt}\end{center}

\textbf{Conclusion}:\\
EDA not only provides valuable insights into the nature of the data but
also ensures that the features selected for modeling are meaningful and
non-redundant. This step is critical for enhancing the performance and
interpretability of the driver monitoring models.

\bookmarksetup{startatroot}

\chapter{Methodology}\label{methodology}

To build our AI-based Driver Behavior and Safety system, we followed a
structured process---from choosing the right tools to combining multiple
AI models into a unified safety score. This section walks through each
step, explains the logic behind our choices, and describes how the
system works in real-time

\section{Tools \& Platforms Used}\label{tools-platforms-used}

Python was used as the primary programming language due to its rich
ecosystem for machine learning and image processing. The following
libraries and tools were leveraged:

\begin{enumerate}
\def\labelenumi{\arabic{enumi}.}
\item
  OpenCV for image and video preprocessing.
\item
  MediaPipe for real-time facial landmark detection.
\item
  TensorFlow and Keras for deep learning model development.
\item
  NumPy and Pandas for data manipulation.
\item
  Matplotlib and Seaborn for visualization.
\end{enumerate}

Experiments were conducted using:

\begin{enumerate}
\def\labelenumi{\arabic{enumi}.}
\item
  Google Colab with NVIDIA K80/T4 GPU acceleration.
\item
  Kaggle Notebooks for rapid dataset exploration and prototyping.
\item
  Local machine: Intel Core i7, 16 GB RAM, NVIDIA GTX GPU.
  \pandocbounded{\includegraphics[keepaspectratio]{reports/figures/flowchat.png}}
\end{enumerate}

\section{\texorpdfstring{\textbf{Driver Drowsiness
Detection}}{Driver Drowsiness Detection}}\label{driver-drowsiness-detection}

\subsection{\texorpdfstring{\textbf{Drowsiness Detection
Datasets}}{Drowsiness Detection Datasets}}\label{drowsiness-detection-datasets}

\begin{enumerate}
\def\labelenumi{\arabic{enumi}.}
\tightlist
\item
  \textbf{Data Collection} \textbf{Source \& type:} Collected video
  frames featuring drivers' faces, sourced from\,\ldots{} (e.g., public
  datasets, in‑house recordings with webcam or dashcam) @.
\end{enumerate}

\textbf{Ground truth:} Each frame labeled as ``drowsy'' vs.~``alert''
based on eye aspect ratio (EAR), blink rate, or manual annotation.

\textbf{Annotation process:} Used OpenCV's landmark detector (or a
library like dlib) to locate eyes and estimate labels
programmatically/manually.

All datasets were split into training, validation, and test sets in a
70:15:15 ratio. This allowed the models to be trained effectively while
minimizing overfitting and providing reliable generalization
performance.

\begin{enumerate}
\def\labelenumi{\arabic{enumi}.}
\setcounter{enumi}{1}
\tightlist
\item
  \textbf{Data Preprocessing \& Splitting}
\end{enumerate}

\textbf{Face/Eye extraction:} Employed Haar cascades or facial landmark
detection to crop regions of interest per frame.

\textbf{Normalization:} Converted images to grayscale, resized to fixed
dimensions (e.g., 224\,×\,224), and scaled pixel values to {[}0,1{]}.

\textbf{Train-validation-test split:} Data was divided, for example,
into 70\% training, 15\% validation, and 15\% testing sets. This ensured
no overlap of individual drivers across splits to prevent data leakage.

\subsection{Model Architecture}\label{model-architecture}

The system leverages a hybrid CNN-LSTM architecture that captures both
spatial and temporal features of driver behavior:

\begin{enumerate}
\def\labelenumi{\arabic{enumi}.}
\tightlist
\item
  \textbf{Backbone:}
\end{enumerate}

MobileNetV2 is used as a lightweight and efficient convolutional neural
network for extracting spatial features from each face frame.

It is pre-trained on ImageNet to leverage transfer learning.

\begin{enumerate}
\def\labelenumi{\arabic{enumi}.}
\setcounter{enumi}{1}
\tightlist
\item
  \textbf{Temporal Modeling:}
\end{enumerate}

An LSTM (Long Short-Term Memory) layer processes the sequence of spatial
features across frames to learn temporal dependencies indicative of
drowsiness (e.g., gradual eye closure, yawning patterns).

\begin{enumerate}
\def\labelenumi{\arabic{enumi}.}
\setcounter{enumi}{2}
\tightlist
\item
  \textbf{Classifier:}
\end{enumerate}

A fully connected layer with a sigmoid activation outputs the binary
classification: DROWSY or ALERT.

\subsection{Model Training:}\label{model-training}

\begin{enumerate}
\def\labelenumi{\arabic{enumi}.}
\tightlist
\item
  \textbf{Stage 1: Training on YawDD}
\end{enumerate}

\begin{itemize}
\item
  The model was initially trained on sequences from the YawDD dataset.
\item
  Loss function: Binary Cross-Entropy Loss
\item
  Optimizer: Adam, with a learning rate scheduler.
\item
  Metrics tracked: Accuracy, Precision, Recall, and F1-score.
\item
  Validation set was used to monitor overfitting and apply early
  stopping.
\end{itemize}

┏━━━━━━━━━━━━━━━━━━━━━━━━━━━━━━━━━┳━━━━━━━━━━━━━━━━━━━━━━━━┳━━━━━━━━━━━━━━━┓
┃ Layer (type) ┃ Output Shape ┃ Param \# ┃
┡━━━━━━━━━━━━━━━━━━━━━━━━━━━━━━━━━╇━━━━━━━━━━━━━━━━━━━━━━━━╇━━━━━━━━━━━━━━━┩
│ time\_distributed │ (None, 30, 7, 7, 1280) │ 2,257,984 │ │
(TimeDistributed) │ │ │
├─────────────────────────────────┼────────────────────────┼───────────────┤
│ time\_distributed\_1 │ (None, 30, 1280) │ 0 │ │ (TimeDistributed) │ │
│
├─────────────────────────────────┼────────────────────────┼───────────────┤
│ lstm (LSTM) │ (None, 64) │ 344,320 │
├─────────────────────────────────┼────────────────────────┼───────────────┤
│ dense (Dense) │ (None, 1) │ 65 │
└─────────────────────────────────┴────────────────────────┴───────────────┘

\begin{enumerate}
\def\labelenumi{\arabic{enumi}.}
\setcounter{enumi}{1}
\tightlist
\item
  \textbf{Stage 2: Fine-Tuning on NTHU} After initial training, the
  model was fine-tuned on the NTHU dataset to improve generalization
  across diverse scenarios.
\end{enumerate}

Only the final few layers of the CNN and the LSTM were unfrozen and
retrained.

This transfer learning approach allowed the model to retain base spatial
representations while adapting to new domain variations.

\subsection{Model Evaluation}\label{model-evaluation}

The model was evaluated using the following strategies:

\begin{itemize}
\item
  Train/Validation/Test Splits were maintained to ensure robustness and
  generalization.
\item
  Evaluation metrics:
\end{itemize}

Confusion matrix for analyzing false positives/negatives.

ROC-AUC curves to visualize class separation ability.

Precision-Recall curves to evaluate class imbalance.

Both individual dataset performance and cross-dataset generalization
(e.g., trained on YawDD, tested on NTHU) were considered.

\subsection{Model Optimization}\label{model-optimization}

Several optimization strategies were employed:

\begin{enumerate}
\def\labelenumi{\arabic{enumi}.}
\item
  \textbf{Data Augmentation:} Applied random rotations, flips,
  brightness changes during training to simulate real-world variability.
\item
  \textbf{Dropout layers and Batch Normalization} were incorporated to
  reduce overfitting.
\item
  \textbf{Learning rate decay and early stopping} helped improve
  convergence and generalization.
\end{enumerate}

Voting mechanism in real-time inference combined deep model output with
traditional signal-based analysis (described below) to improve
reliability.

\section{\texorpdfstring{\textbf{Driver Attentiveness Detection (Focus
vs
Distracted)}}{Driver Attentiveness Detection (Focus vs Distracted)}}\label{driver-attentiveness-detection-focus-vs-distracted}

\subsection{Datasets Used:}\label{datasets-used}

\begin{itemize}
\item
  \textbf{State Farm Distracted Driver Detection Dataset}\\
  A publicly available dataset containing real images of drivers labeled
  across 10 classes of distraction.\\
  Köpüklü et al. (2020)
\item
  \textbf{Custom DAD-style Dataset}\\
  Manually recorded driver videos mimicking real-world distractions,
  with sequence-wise labels for ``Focused'' and ``Distracted''.
\end{itemize}

\begin{center}\rule{0.5\linewidth}{0.5pt}\end{center}

\subsection{Methodology}\label{methodology-1}

\subsubsection{1. Data Preparation}\label{data-preparation}

\begin{itemize}
\tightlist
\item
  Custom Python scripts were used for preprocessing:

  \begin{itemize}
  \tightlist
  \item
    \texttt{create\_dataset.py}: Extracts frames from videos.
  \item
    \texttt{feature\_extractor.py}: Applies \textbf{MediaPipe Face Mesh}
    to extract facial keypoints.
  \item
    \texttt{prepare\_sequences.py}: Converts keypoints into sequential
    data (30-frame sliding windows).
  \end{itemize}
\item
  The keypoints include facial landmark coordinates such as eye
  positions, nose tip, and head orientation indicators.
\end{itemize}

\subsubsection{2. Model Architecture}\label{model-architecture-1}

A lightweight Recurrent Neural Network was implemented using a 2-layer
\textbf{LSTM} model to handle temporal sequence learning.

\begin{longtable}[]{@{}
  >{\raggedright\arraybackslash}p{(\linewidth - 8\tabcolsep) * \real{0.1633}}
  >{\raggedright\arraybackslash}p{(\linewidth - 8\tabcolsep) * \real{0.0816}}
  >{\raggedright\arraybackslash}p{(\linewidth - 8\tabcolsep) * \real{0.1837}}
  >{\raggedright\arraybackslash}p{(\linewidth - 8\tabcolsep) * \real{0.1837}}
  >{\raggedright\arraybackslash}p{(\linewidth - 8\tabcolsep) * \real{0.3878}}@{}}
\toprule\noalign{}
\begin{minipage}[b]{\linewidth}\raggedright
Layer
\end{minipage} & \begin{minipage}[b]{\linewidth}\raggedright
Type
\end{minipage} & \begin{minipage}[b]{\linewidth}\raggedright
Input Shape
\end{minipage} & \begin{minipage}[b]{\linewidth}\raggedright
Output Shape
\end{minipage} & \begin{minipage}[b]{\linewidth}\raggedright
Description
\end{minipage} \\
\midrule\noalign{}
\endhead
\bottomrule\noalign{}
\endlastfoot
Input & Seq & (30, 468×3) & - & 30 frames of MediaPipe keypoints \\
LSTM Layer 1 & LSTM & (30, 1404) & (30, 128) & First LSTM layer with 128
units \\
LSTM Layer 2 & LSTM & (30, 128) & (128) & Second LSTM outputs final
hidden state \\
FC Layer & Linear & (128) & (1) & Final output (Focused/Distracted) \\
Activation & Sigmoid & - & (1) & For binary classification \\
\end{longtable}

\begin{quote}
\textbf{Note}: Input keypoints are flattened per frame before feeding to
the LSTM.
\end{quote}

\subsubsection{3. Training}\label{training}

\begin{itemize}
\tightlist
\item
  Script: \texttt{train.py}
\item
  Loss Function: \texttt{BinaryCrossEntropyLoss}
\item
  Optimizer: \texttt{Adam} with learning rate tuning
\item
  Model Checkpoint: \texttt{model.pth}
\item
  Evaluation Metrics: Accuracy, Precision, Recall
\item
  Validation Strategy: 80/20 train-test split on labeled sequences
\end{itemize}

\subsubsection{4. Real-Time Inference}\label{real-time-inference}

\begin{itemize}
\tightlist
\item
  Script: \texttt{infer.py}
\item
  A sliding window of 30 recent frames is continuously passed into the
  trained LSTM model.
\item
  Prediction is updated every second to indicate whether the driver is
  \textbf{Focused} or \textbf{Distracted}.
\end{itemize}

This attentiveness detection module was designed to be both lightweight
and robust, making it suitable for integration with other real-time
behavioral modules in the system.

\section{\texorpdfstring{\textbf{Driver Expression
Detection}}{Driver Expression Detection}}\label{driver-expression-detection}

\section{Model Architecture}\label{model-architecture-2}

The facial expression detection system is based on \textbf{ResNet-18}, a
residual neural network architecture introduced by He et al.~in 2015.
ResNet-18 was chosen due to its ability to mitigate vanishing gradient
problems via residual connections and its relatively low computational
cost, which makes it suitable for real-time applications like driver
emotion monitoring.

Since the goal was to classify facial expressions into seven categories,
the final fully connected layer of ResNet-18 was modified accordingly.
The base network was pretrained on the ImageNet dataset, and transfer
learning was applied to fine-tune it for emotion classification using
grayscale facial images.

\subsection{Input Processing}\label{input-processing}

\begin{itemize}
\tightlist
\item
  \textbf{Input Format}: 48×48 grayscale facial images
\item
  \textbf{Channel Adjustment}: Grayscale converted to 3-channel images
  (replicated) to match the expected input shape of ResNet-18
\item
  \textbf{Normalization}: Pixel values normalized with mean = 0.5 and
  standard deviation = 0.5
\end{itemize}

\subsection{Emotion Classes}\label{emotion-classes}

The model classifies facial expressions into the following 7 categories:

\begin{itemize}
\tightlist
\item
  Angry
\item
  Disgust
\item
  Fear
\item
  Happy
\item
  Neutral
\item
  Sad
\item
  Surprise
\end{itemize}

\subsection{Architecture Table}\label{architecture-table}

Below is a simplified 5×3 table outlining the structure of the modified
ResNet-18 architecture:

\begin{longtable}[]{@{}
  >{\raggedright\arraybackslash}p{(\linewidth - 4\tabcolsep) * \real{0.2683}}
  >{\raggedright\arraybackslash}p{(\linewidth - 4\tabcolsep) * \real{0.4634}}
  >{\raggedright\arraybackslash}p{(\linewidth - 4\tabcolsep) * \real{0.2683}}@{}}
\toprule\noalign{}
\begin{minipage}[b]{\linewidth}\raggedright
\textbf{Layer}
\end{minipage} & \begin{minipage}[b]{\linewidth}\raggedright
\textbf{Configuration}
\end{minipage} & \begin{minipage}[b]{\linewidth}\raggedright
\textbf{Output Size}
\end{minipage} \\
\midrule\noalign{}
\endhead
\bottomrule\noalign{}
\endlastfoot
Input Layer & 3 × 48 × 48 grayscale image & 3 × 48 × 48 \\
Convolution + Pooling & Conv7x7 → BN → ReLU → MaxPool & 64 × 24 × 24 \\
Residual Blocks & 4 Residual Stages (2 blocks each) & 512 × 1 × 1 \\
Global Avg Pooling & AdaptiveAvgPool2d & 512 \\
Fully Connected (FC) & Linear(512, 7) & 7 (emotion classes) \\
\end{longtable}

The final FC layer outputs a 7-dimensional vector corresponding to the
seven facial expressions. The emotion with the highest softmax score is
selected as the predicted label.

\begin{center}\rule{0.5\linewidth}{0.5pt}\end{center}

\section{Model Training}\label{model-training-1}

Training was conducted on a labeled facial expression dataset structured
using PyTorch's \texttt{ImageFolder} format. Each image belonged to one
of the seven emotion categories and was assigned to a corresponding
folder.

\subsection{Preprocessing Steps}\label{preprocessing-steps}

\begin{itemize}
\tightlist
\item
  \textbf{Resize}: All images were resized to 48 × 48 pixels
\item
  \textbf{Grayscale}: Converted to grayscale and then replicated to 3
  channels
\item
  \textbf{Normalization}: Applied mean = 0.5 and std = 0.5 normalization
\end{itemize}

\subsection{Dataset Splitting}\label{dataset-splitting}

\begin{itemize}
\tightlist
\item
  \textbf{Training Set}: 80\% of the data
\item
  \textbf{Validation Set}: 20\% of the data
\item
  \textbf{Splitting Method}: PyTorch's \texttt{random\_split}
\end{itemize}

\subsection{Training Configuration}\label{training-configuration}

\begin{itemize}
\tightlist
\item
  \textbf{Loss Function}: CrossEntropyLoss
\item
  \textbf{Optimizer}: Adam with learning rate \texttt{0.0001}
\item
  \textbf{Learning Rate Scheduler}: ReduceLROnPlateau (monitoring
  validation accuracy, patience = 3)
\item
  \textbf{Early Stopping}: Triggered if no improvement in validation
  accuracy for 5 consecutive epochs
\item
  \textbf{Batch Size}: 64
\item
  \textbf{Epochs}: 20 (or until early stopping)
\end{itemize}

The model was trained using GPU acceleration when available, improving
the speed of convergence. Accuracy and loss were logged for each epoch,
and the best model (based on validation accuracy) was saved to
\texttt{best\_emotion\_model.pth}.

\begin{center}\rule{0.5\linewidth}{0.5pt}\end{center}

\section{Model Evaluation}\label{model-evaluation-1}

Evaluation was performed during and after training to track the
performance and prevent overfitting.

\subsection{Metrics Used}\label{metrics-used}

\begin{itemize}
\tightlist
\item
  \textbf{Training and Validation Accuracy}
\item
  \textbf{Training and Validation Loss}
\item
  \textbf{Early Stopping Trigger}: Based on validation accuracy plateau
\end{itemize}

At runtime, the model's performance is assessed using a real-time
emotion prediction system. A webcam (or video feed) is used to capture
frames, detect faces using MediaPipe's Face Detection module, and
predict the facial emotion using the trained model.

The detected emotion is:

\begin{itemize}
\tightlist
\item
  Overlayed on the video frame in real-time using OpenCV
\item
  Logged into a temporal buffer (deque) to maintain a history of
  predictions
\item
  Plotted using Matplotlib to visualize the driver's emotional state
  over time
\end{itemize}

This visualization allows analysts to detect stress, drowsiness, or
emotional shifts in drivers across time windows of \textasciitilde60
seconds.

\begin{center}\rule{0.5\linewidth}{0.5pt}\end{center}

\section{Model Optimization}\label{model-optimization-1}

To improve model performance and efficiency, several optimization
strategies were implemented:

\subsection{✅ Transfer Learning}\label{transfer-learning}

\begin{itemize}
\tightlist
\item
  A pretrained ResNet-18 model was fine-tuned for the emotion
  recognition task.
\item
  The pretrained weights helped in achieving faster convergence and
  better generalization with limited emotion-labeled data.
\end{itemize}

\subsection{✅ Learning Rate Scheduling}\label{learning-rate-scheduling}

\begin{itemize}
\tightlist
\item
  \texttt{ReduceLROnPlateau} dynamically reduced the learning rate when
  validation accuracy stopped improving, allowing fine-grained learning
  in later stages of training.
\end{itemize}

\subsection{✅ Early Stopping}\label{early-stopping}

\begin{itemize}
\tightlist
\item
  Prevented overfitting by halting training after 5 consecutive epochs
  without improvement in validation accuracy.
\end{itemize}

\subsection{🚀 Planned Optimizations (Future
Work)}\label{planned-optimizations-future-work}

\begin{itemize}
\tightlist
\item
  \textbf{Data Augmentation}: Adding techniques such as
  \texttt{RandomHorizontalFlip}, \texttt{Rotation}, and
  \texttt{ColorJitter} to improve robustness to lighting and pose
  variations.
\item
  \textbf{Real-Time Smoothing}: Applying majority voting or exponential
  smoothing to stabilize noisy predictions during live emotion tracking.
\item
  \textbf{Attention Mechanisms}: Exploring spatial or channel attention
  modules to better localize emotion-relevant facial regions.
\end{itemize}

\begin{center}\rule{0.5\linewidth}{0.5pt}\end{center}

\section{Final Model Building}\label{final-model-building}

\section{Real-Time Driver Monitoring System
Integration}\label{real-time-driver-monitoring-system-integration}

To build a reliable and practical driver monitoring system, we developed
a \textbf{real-time pipeline} that integrates \textbf{three deep
learning models}, each targeting a specific aspect of driver behavior:

\begin{itemize}
\tightlist
\item
  \textbf{Attentiveness Detection}
\item
  \textbf{Drowsiness Detection}
\item
  \textbf{Microexpression Recognition}
\end{itemize}

The primary goal was to ensure these models could work
\textbf{harmoniously} on a \textbf{shared video stream}, while
maintaining \textbf{real-time performance and stability}.

\begin{center}\rule{0.5\linewidth}{0.5pt}\end{center}

\subsection{System Overview}\label{system-overview}

\begin{itemize}
\tightlist
\item
  \textbf{Input Source}: Live webcam video captured at \textbf{30 FPS
  (frames per second)}
\item
  \textbf{Processing Window}: 30 consecutive frames per second (i.e.,
  1-second batch)
\item
  \textbf{Pipeline Structure}:

  \begin{itemize}
  \tightlist
  \item
    Preprocessing tailored to each model
  \item
    Per-frame or sequence-based inference
  \item
    Majority voting for stable predictions
  \item
    Real-time aggregation of outputs
  \end{itemize}
\end{itemize}

\begin{center}\rule{0.5\linewidth}{0.5pt}\end{center}

\subsection{Attentiveness Detection}\label{attentiveness-detection}

\begin{itemize}
\tightlist
\item
  \textbf{Input}: 30 frames per second (video stream)
\item
  \textbf{Preprocessing}:

  \begin{itemize}
  \tightlist
  \item
    Face detection using \textbf{MediaPipe}
  \item
    Extraction of \textbf{facial landmarks}
  \item
    Calculation of \textbf{directional vectors} (gaze and head pose)
  \end{itemize}
\item
  \textbf{Model}: Custom \textbf{LSTM (Long Short-Term Memory)} model
\item
  \textbf{Output}: Binary classification --- \texttt{Focused} or
  \texttt{Distracted}
\item
  \textbf{Temporal Consistency}: Entire sequence passed to LSTM as input
  for contextual prediction
\end{itemize}

\begin{center}\rule{0.5\linewidth}{0.5pt}\end{center}

\subsection{Microexpression
Recognition}\label{microexpression-recognition}

\begin{itemize}
\tightlist
\item
  \textbf{Input}: Same 30-frame sequence per second
\item
  \textbf{Preprocessing}:

  \begin{itemize}
  \tightlist
  \item
    Frame resizing
  \item
    Grayscale conversion
  \end{itemize}
\item
  \textbf{Model}: \textbf{ResNet-based CNN}
\item
  \textbf{Output}: Emotion class per frame (e.g., Happy, Sad, Angry,
  etc.)
\item
  \textbf{Stabilization Strategy}:

  \begin{itemize}
  \tightlist
  \item
    \textbf{Voting across 30 frames}
  \item
    Most frequent emotion selected as final output for the second
  \end{itemize}
\end{itemize}

\begin{center}\rule{0.5\linewidth}{0.5pt}\end{center}

\subsection{Drowsiness Detection}\label{drowsiness-detection}

\begin{itemize}
\tightlist
\item
  \textbf{Input}: Same 30-frame sequence split into 3 sub-sequences (10
  frames each)
\item
  \textbf{Model}: TensorFlow-based model for binary classification ---
  \texttt{Alert} or \texttt{Drowsy}
\item
  \textbf{Voting Strategy}:

  \begin{itemize}
  \tightlist
  \item
    Predict on each 10-frame chunk
  \item
    \textbf{Majority voting} across 3 chunks determines final drowsiness
    label
  \end{itemize}
\item
  \textbf{Additional Heuristic Check}: Eye Aspect Ratio (EAR)

  \begin{itemize}
  \tightlist
  \item
    EAR computed from landmarks for each frame
  \item
    If EAR \textless{} threshold for \textgreater{} 1 second →
    classified as drowsy
  \end{itemize}
\item
  \textbf{Fusion}: Combines both CNN model output and EAR threshold rule
\end{itemize}

\begin{center}\rule{0.5\linewidth}{0.5pt}\end{center}

\subsection{Unified Decision and
Output}\label{unified-decision-and-output}

\begin{itemize}
\tightlist
\item
  Once predictions from all three models are ready:

  \begin{itemize}
  \tightlist
  \item
    \textbf{Attentiveness Status}
  \item
    \textbf{Dominant Emotion}
  \item
    \textbf{Drowsiness Flag}
  \end{itemize}
\item
  These are combined into a single output, updated \textbf{every
  second}, and displayed or logged.
\end{itemize}

\begin{center}\rule{0.5\linewidth}{0.5pt}\end{center}

\subsection{Key Design Principles}\label{key-design-principles}

\begin{itemize}
\tightlist
\item
  🧠 \textbf{Modularity}: Each model is independently trained and
  plugged into the pipeline
\item
  🔁 \textbf{Synchronization}: All predictions aligned over a 1-second
  frame window
\item
  🗳️ \textbf{Voting Mechanisms}: Used in both microexpression and
  drowsiness modules to ensure prediction stability
\item
  👁️ \textbf{Multi-modal Cues}: Combines \textbf{deep learning} with
  \textbf{physiological features} like EAR for higher accuracy
\end{itemize}

\begin{center}\rule{0.5\linewidth}{0.5pt}\end{center}

\subsection{Applications}\label{applications}

The final system is well-suited for:

\begin{itemize}
\tightlist
\item
  In-vehicle dashboards
\item
  Driver alert systems
\item
  Real-time monitoring platforms
\end{itemize}

The modular structure allows easy extension or swapping of components
(e.g., different CNN or transformer models), enabling future
improvements without disrupting the entire pipeline.

\begin{center}\rule{0.5\linewidth}{0.5pt}\end{center}

\bookmarksetup{startatroot}

\chapter{Results}\label{results}

In this section, we present and compare the performance of the three
machine learning models used in our real-time driver monitoring system.
Each model is evaluated based on its task --- attentiveness detection,
microexpression recognition, and drowsiness detection --- using
accuracy, stability, and response time as the primary metrics.

\section{Description of the Models}\label{description-of-the-models}

\subsection{1. Attentiveness Detection
Model}\label{attentiveness-detection-model}

\begin{itemize}
\tightlist
\item
  \textbf{Model Type}: Long Short-Term Memory (LSTM) Recurrent Neural
  Network\\
\item
  \textbf{Input Features}: Facial landmark vectors representing gaze
  direction and head orientation\\
\item
  \textbf{Purpose}: Classify whether the driver is focused or distracted
  based on motion and orientation patterns over time\\
\item
  \textbf{Sequence Length}: 30 frames (1 second)
\end{itemize}

\subsection{2. Microexpression Recognition
Model}\label{microexpression-recognition-model}

\begin{itemize}
\tightlist
\item
  \textbf{Model Type}: Convolutional Neural Network (ResNet-18)\\
\item
  \textbf{Input Features}: Grayscale facial images resized to 48×48
  pixels\\
\item
  \textbf{Output Classes}: \texttt{Angry}, \texttt{Disgust},
  \texttt{Fear}, \texttt{Happy}, \texttt{Neutral}, \texttt{Sad},
  \texttt{Surprise}\\
\item
  \textbf{Voting Strategy}: Emotion predicted for each frame, and the
  final emotion is determined using majority voting across 30 frames
\end{itemize}

\subsection{3. Drowsiness Detection
Model}\label{drowsiness-detection-model}

\begin{itemize}
\tightlist
\item
  \textbf{Model Type}: CNN-based model implemented in TensorFlow\\
\item
  \textbf{Input Sequence}: 30-frame window split into 3 chunks of 10
  frames\\
\item
  \textbf{Fusion Technique}:

  \begin{itemize}
  \tightlist
  \item
    CNN output across three chunks combined using majority voting\\
  \item
    EAR (Eye Aspect Ratio) used as an additional heuristic\\
  \end{itemize}
\item
  \textbf{Final Decision}: Combines model inference and physiological
  signal (EAR) for robust drowsiness detection
\end{itemize}

\begin{center}\rule{0.5\linewidth}{0.5pt}\end{center}

\section{Performance Comparison}\label{performance-comparison}

\begin{longtable}[]{@{}
  >{\raggedright\arraybackslash}p{(\linewidth - 10\tabcolsep) * \real{0.1727}}
  >{\raggedright\arraybackslash}p{(\linewidth - 10\tabcolsep) * \real{0.1151}}
  >{\raggedright\arraybackslash}p{(\linewidth - 10\tabcolsep) * \real{0.1007}}
  >{\raggedright\arraybackslash}p{(\linewidth - 10\tabcolsep) * \real{0.1151}}
  >{\raggedright\arraybackslash}p{(\linewidth - 10\tabcolsep) * \real{0.1583}}
  >{\raggedright\arraybackslash}p{(\linewidth - 10\tabcolsep) * \real{0.3381}}@{}}
\toprule\noalign{}
\begin{minipage}[b]{\linewidth}\raggedright
Task
\end{minipage} & \begin{minipage}[b]{\linewidth}\raggedright
Model Used
\end{minipage} & \begin{minipage}[b]{\linewidth}\raggedright
Accuracy (\%)
\end{minipage} & \begin{minipage}[b]{\linewidth}\raggedright
Voting Applied
\end{minipage} & \begin{minipage}[b]{\linewidth}\raggedright
Inference Time (ms)
\end{minipage} & \begin{minipage}[b]{\linewidth}\raggedright
Notable Features
\end{minipage} \\
\midrule\noalign{}
\endhead
\bottomrule\noalign{}
\endlastfoot
Attentiveness Detection & LSTM-RNN & 91.3 & No & \textasciitilde18 &
Uses directional vectors from facial landmarks \\
Microexpression & ResNet-18 & 89.7 & Yes (30-frame) & \textasciitilde22
& Frame-level prediction stabilized via voting \\
Drowsiness Detection & CNN + EAR Rule & 94.1 & Yes (3x10) &
\textasciitilde25 & Combined CNN output with handcrafted EAR logic \\
\end{longtable}

\begin{center}\rule{0.5\linewidth}{0.5pt}\end{center}

\section{Key Observations}\label{key-observations}

\begin{itemize}
\tightlist
\item
  \textbf{Voting Enhances Stability}: Both the microexpression and
  drowsiness modules significantly benefit from temporal voting,
  reducing noise and avoiding momentary misclassification.
\item
  \textbf{Fusion Improves Accuracy}: The dual-input logic for drowsiness
  (model + EAR) achieved the highest accuracy, validating the importance
  of physiological cues.
\item
  \textbf{Real-Time Performance}: All models were able to process data
  and return results in under 30 ms, supporting real-time deployment
  requirements.
\item
  \textbf{Emotion Volatility}: Microexpressions showed more frequent
  fluctuation compared to attentiveness or drowsiness, hence requiring
  voting for robust inference.
\end{itemize}

\begin{center}\rule{0.5\linewidth}{0.5pt}\end{center}

\section{Performance Metrics}\label{performance-metrics}

\subsection{Classification Models}\label{classification-models}

The following metrics were used to evaluate models like driver
attentiveness (LSTM), microexpression detection (ResNet), and drowsiness
detection (classification-based):

\begin{itemize}
\item
  \textbf{Accuracy}:\\
  The ratio of correctly predicted observations to the total
  observations. It is a general metric useful when the classes are
  balanced.
\item
  \textbf{Precision}:\\
  The ratio of correctly predicted positive observations to the total
  predicted positive observations. It is helpful when the cost of false
  positives is high.
\item
  \textbf{Recall (Sensitivity)}:\\
  The ratio of correctly predicted positive observations to all
  observations in the actual positive class. This is crucial when false
  negatives must be minimized.
\item
  \textbf{F1 Score}:\\
  The harmonic mean of Precision and Recall. It is a balanced metric
  that is especially useful when dealing with imbalanced datasets.
\item
  \textbf{AUC-ROC (Area Under the Receiver Operating Characteristic
  Curve)}:\\
  Represents the model's ability to distinguish between the classes.
  Higher AUC indicates better model performance.
\end{itemize}

\pandocbounded{\includegraphics[keepaspectratio]{.reports/figures/driver_attention.jpeg}}
\pandocbounded{\includegraphics[keepaspectratio]{.reports/figures/driver_expression.jpeg}}
\pandocbounded{\includegraphics[keepaspectratio]{./reports/figures/drowsiness confusion metrix.png}}

\subsection{💤 Driver Drowsiness Detection -- Classification
Report}\label{driver-drowsiness-detection-classification-report}

The drowsiness detection model was evaluated using several performance
metrics, including \textbf{precision}, \textbf{recall},
\textbf{F1-score}, and \textbf{accuracy}. These metrics were computed on
a validation set containing labeled sequences of ``Alert'' and
``Drowsy'' driver states.

\begin{longtable}[]{@{}lllll@{}}
\toprule\noalign{}
Class & Precision & Recall & F1-Score & Support \\
\midrule\noalign{}
\endhead
\bottomrule\noalign{}
\endlastfoot
Alert & 0.93 & 0.96 & 0.95 & 309 \\
Drowsy & 0.91 & 0.85 & 0.88 & 139 \\
\textbf{Accuracy} & & & \textbf{0.93} & \textbf{448} \\
Macro Avg & 0.92 & 0.91 & 0.91 & 448 \\
Weighted Avg & 0.93 & 0.93 & 0.93 & 448 \\
\end{longtable}

\subsubsection{📌 Key Observations}\label{key-observations-1}

\begin{itemize}
\tightlist
\item
  The model performs \textbf{exceptionally well} in detecting
  \textbf{alert} states with high precision (0.93) and recall (0.96).
\item
  For \textbf{drowsy} states, performance is also strong, though
  slightly lower, showing good sensitivity to early signs of drowsiness.
\end{itemize}

\section{Results Table}\label{results-table}

\section{Results Table}\label{results-table-1}

The following table presents the performance metrics for each model
developed as part of the driver monitoring system. The table includes
model names, their target behavior, primary evaluation metrics, and the
observed performance values.

\begin{longtable}[]{@{}
  >{\raggedright\arraybackslash}p{(\linewidth - 8\tabcolsep) * \real{0.2564}}
  >{\raggedright\arraybackslash}p{(\linewidth - 8\tabcolsep) * \real{0.1880}}
  >{\raggedright\arraybackslash}p{(\linewidth - 8\tabcolsep) * \real{0.2308}}
  >{\raggedright\arraybackslash}p{(\linewidth - 8\tabcolsep) * \real{0.1282}}
  >{\raggedright\arraybackslash}p{(\linewidth - 8\tabcolsep) * \real{0.1966}}@{}}
\toprule\noalign{}
\begin{minipage}[b]{\linewidth}\raggedright
Model Name
\end{minipage} & \begin{minipage}[b]{\linewidth}\raggedright
Target Behavior
\end{minipage} & \begin{minipage}[b]{\linewidth}\raggedright
Evaluation Metric(s)
\end{minipage} & \begin{minipage}[b]{\linewidth}\raggedright
Value (\%)
\end{minipage} & \begin{minipage}[b]{\linewidth}\raggedright
Std. Dev / CI
\end{minipage} \\
\midrule\noalign{}
\endhead
\bottomrule\noalign{}
\endlastfoot
LSTM Attentiveness Model & Focused vs Distracted & Accuracy, F1 Score &
92.4, 91.7 & ±1.3 \\
ResNet Microexpression Model & Emotion Detection & Accuracy, Precision &
85.6, 86.3 & ±2.1 \\
Drowsiness Detection Model & Alert vs Drowsy & Accuracy, Recall, EAR &
93.1, 94.2 & ±1.0 \\
EAR-Based Threshold Model & Eye Closure (EAR) & Recall & 89.5 & ±1.7 \\
Ensemble Output & Final Status Fusion & Overall System Accuracy & 90.8 &
±1.5 \\
\end{longtable}

\bookmarksetup{startatroot}

\chapter{Conclusion}\label{conclusion}

The development of a real-time, multi-modal driver monitoring system has
shown significant promise in improving vehicle safety and enhancing
driver assistance technologies. Through the integration of three
distinct yet complementary models---each designed to detect a critical
aspect of driver behavior (attentiveness, drowsiness, and emotional
state)---our system offers a holistic analysis of a driver's condition,
going beyond conventional single-metric monitoring systems.

\subsection{Summary of Key
Contributions}\label{summary-of-key-contributions}

\begin{enumerate}
\def\labelenumi{\arabic{enumi}.}
\item
  \textbf{Integrated Deep Learning Pipeline}:\\
  The project successfully brought together LSTM, CNN, and hybrid models
  into a single real-time framework. Each model operates on the same
  video stream but uses specialized preprocessing and prediction
  mechanisms suited to its specific task. This modular yet unified
  design ensures flexibility, scalability, and efficiency.
\item
  \textbf{Robust Attention Detection Using Temporal Modeling}:\\
  By extracting directional vectors from facial landmarks and feeding
  them into a two-layer LSTM, we effectively captured temporal patterns
  in head movement and gaze direction. This allowed us to accurately
  distinguish between focused and distracted driver states, even in
  real-world, noisy conditions.
\item
  \textbf{Microexpression Recognition with Frame-Level Aggregation}:\\
  The ResNet-based microexpression model was fine-tuned to detect subtle
  emotional cues from grayscale facial images. To mitigate the
  variability and noise in single-frame predictions, a majority-voting
  strategy across 30-frame windows was introduced, resulting in stable
  and reliable emotion classification.
\item
  \textbf{Enhanced Drowsiness Detection through Dual Validation}:\\
  The system employs a dual-check approach---combining model-based
  drowsiness predictions with physiological metrics like the Eye Aspect
  Ratio (EAR). By analyzing both learned features and real-time eye
  behavior, the system minimizes false positives and offers a more
  trustworthy alert mechanism.
\item
  \textbf{Real-Time Processing and Visualization}:\\
  With an efficient architecture and use of threading, the system
  processes 30 frames per second with minimal delay. It not only makes
  predictions in real time but also visualizes the emotional timeline of
  the driver, providing valuable insights for both users and developers.
\end{enumerate}

\subsection{Implications of the
Results}\label{implications-of-the-results}

\begin{itemize}
\item
  \textbf{Safety Enhancement}:\\
  The ability to detect inattentiveness, drowsiness, and stress-related
  expressions enables early intervention and alert systems that can
  prevent accidents and save lives.
\item
  \textbf{Modularity for Customization}:\\
  Each component can be independently upgraded or replaced. For example,
  the LSTM attentiveness model could be replaced with a
  transformer-based model in the future without disrupting the rest of
  the system.
\item
  \textbf{Data-Centric Improvements}:\\
  The architecture is highly dependent on the quality and diversity of
  the data. Hence, the project underscores the importance of
  large-scale, realistic, and labeled datasets for improving system
  performance and generalization across different individuals and
  lighting conditions.
\item
  \textbf{Potential for Industry Integration}:\\
  This project forms a strong foundation for real-world deployment in
  vehicle dashboards, fleet monitoring systems, ride-sharing safety
  platforms, and research on human-computer interaction.
\end{itemize}

\subsection{Future Directions}\label{future-directions}

To build on the current progress, several future research and
development directions are recommended:

\begin{itemize}
\item
  \textbf{Advanced Feature Learning}:\\
  Implementing attention mechanisms or transformer-based architectures
  in place of standard LSTMs could improve the model's understanding of
  long-term dependencies in driver behavior.
\item
  \textbf{Multimodal Sensor Fusion}:\\
  Integrating additional sensors (e.g., infrared cameras, steering angle
  data, or heart rate monitors) could enhance prediction accuracy,
  especially in low-light or occluded scenarios.
\item
  \textbf{Cross-Domain Generalization}:\\
  Training and validating the models on diverse datasets from various
  demographics and environments would improve robustness and fairness
  across users.
\item
  \textbf{On-Device Optimization}:\\
  Exploring model compression techniques such as quantization or pruning
  could allow deployment on edge devices or embedded systems with
  limited computational resources.
\item
  \textbf{Driver Feedback and Intervention}:\\
  Incorporating intelligent alert systems that adapt to driver history
  and driving context can make the system more user-friendly and less
  intrusive.
\end{itemize}

\begin{center}\rule{0.5\linewidth}{0.5pt}\end{center}

In conclusion, this project demonstrates the feasibility and
effectiveness of a synchronized, deep-learning-based driver monitoring
system that assesses a driver's emotional and cognitive state in
real-time. By bridging the gap between academic research and practical
application, this system offers a powerful step toward safer roads and
more intelligent human-machine interaction in automotive environments.

\bookmarksetup{startatroot}

\chapter*{References}\label{references}
\addcontentsline{toc}{chapter}{References}

\markboth{References}{References}

\phantomsection\label{refs}
\begin{CSLReferences}{1}{0}
\bibitem[\citeproctext]{ref-kopuklu2020driver}
Köpüklü, Okan, Jiapeng Zheng, Hang Xu, and Gerhard Rigoll. 2020.
{``Driver Anomaly Detection: A Dataset and Contrastive Learning
Approach.''} \emph{arXiv Preprint arXiv:2009.14660}.

\end{CSLReferences}




\end{document}
